\subsubsection{Manifoldlikeness}

\begin{theorem}[Causal manifoldlikeness --- the order-complex cohomology as spacetime topology]\label{thm:causal-manifoldlikeness}\lean{Infomaterial.Relational.Spacetime.Dimension.Manifoldlikeness.causal_manifoldlikeness}\leanok\uses{thm:causal-cohomology-n, thm:mm-ordering-fraction}
General $H^n$ (\cref{thm:causal-cohomology-n}) is the instrument for
manifoldlikeness: in causal-set theory the cohomology of the order complex is
the topology of the reconstructed spacetime. Two halves are proved. First, the
topology is MLD-invariant: pullback along any factor map is a cochain map
(\texttt{pullback\_delta}), hence carries cocycles to cocycles and
coboundaries to coboundaries. Second, the computable invariant is the Euler
characteristic $\sum_k (-1)^k\,\#\{k\text{-chains}\}$, which by Philip Hall is
the incidence M\"obius function (\texttt{hallSum} = \texttt{mobius}). On the
four-event diamond, the interval $(0,3)$ has chain census
$(c_1,c_2,c_3) = (1,2,0)$, so $\chi = +1 = \mu(0,3)$: the order complex of
the interval is the two spacelike midpoints $\{1,2\}$, and $\mu = +1 =
(-1)^2$ is the link signature of a 2-dimensional Alexandrov interval.
\end{theorem}

\begin{theorem}[Causal-set Hauptvermutung --- the continuum-limit Betti match]\label{thm:causal-set-hauptvermutung}\notready\uses{thm:causal-manifoldlikeness}
\emph{Open frontier, stated but not formalized.}
\cref{thm:causal-manifoldlikeness} makes the manifoldlikeness question precise
and proves its invariance and Euler-characteristic halves, but not the
Hauptvermutung itself: that a causal set sprinkled into a Lorentzian manifold
$M$ has, in the large-$N$ limit, order-complex cohomology matching the topology
of $M$, so that a causal set looks like at most one manifold. The missing
pieces are finite-rank causal cohomology, Euler--Poincar\'e over the full
complex, and the probabilistic continuum-limit convergence.
\end{theorem}
